% !TEX root = ../thesis.tex
% packages
\usepackage{amsmath,amsfonts,amssymb}
\usepackage{mathrsfs}
\usepackage{mathtools}
\usepackage{simplebnf}
\usepackage{caption}
\usepackage{mathpartir}
\usepackage{stackengine}
\usepackage[dvipsnames]{xcolor}
\usepackage[shortlabels,inline]{enumitem}

% figure
\usepackage{tikz}
\usepackage{tikz-cd}
\usetikzlibrary{positioning,arrows.meta}
\usepackage{float}
\usepackage{graphicx}
\graphicspath{ {figures/} }

\tikzset{
  logicnode/.style={
    rectangle,
    rounded corners=2pt,
    draw=black,
    fill=gray!8,
    minimum width=1.55cm,
    minimum height=8mm,
    align=center,
    font=\normalsize
  },
  relation/.style={draw=black, line width=0.6pt},
  morphism/.style={-{Stealth[length=2.4mm]}, draw=black, line width=0.7pt},
  dottedmorphism/.style={morphism, dotted, line width=0.9pt},
  edgelabel/.style={font=\normalsize, fill=white, inner sep=1.2pt}
}

% layout
\footskip 4.ex
% ***********************************************************************
\newcommand{\Normalstretch}{1.0} %change to define space between lines
\renewcommand{\baselinestretch}{\Normalstretch}
\newenvironment{singlespace}
{\renewcommand{\baselinestretch}{1} \small \normalsize}
{\renewcommand{\baselinestretch}{\Normalstretch} \small \normalsize}
\renewcommand{\arraystretch}{1.2}
% ***********************************************************************
\newcommand{\baseenvskip}{\baselineskip 2mm}
\usepackage[
  paper=a4paper,
  inner=2.5cm,  % Replaces \oddsidemargin
  outer=2.5cm,  % Replaces \evensidemargin
  top=2.5cm,
  bottom=2.5cm
]{geometry}



% ------------------------------ reference style ----------------------------- %
\usepackage[numbers,sort]{natbib} 	% automatically sort the citations
%\renewcommand\bibname{Bibliography}

% ------------------------------ footer & header ----------------------------- %
% \usepackage{fancyhdr}
% % def headstyles
% \newcommand{\headstyle}{%
% 	\fancyhead[L]{}
% 	\fancyhead[C]{}
% 	\fancyhead[R]{}
% }
% % def footstyles
% \makeatletter
% \newcommand{\footstyle}{%
% 	\fancyfoot[L]{}
% 	\fancyfoot[C]{\thepage}
% 	\fancyfoot[R]{}
% }
% \newcommand{\authorfootstyle}{%
% 	\fancyfoot[L]{}
% 	\fancyfoot[C]{\thepage}
% 	\fancyfoot[R]{}
% }
% \makeatother
% % set global pagestyle
% \pagestyle{fancy}
% \fancyhf{} % clear the orginal headfoot style
% \headstyle
% \footstyle
% % def first-page pagestyle
% \fancypagestyle{firstpage}{%
% 	\fancyhf{}
% 	\headstyle
% 	\authorfootstyle
% }
% % def main pagestyle
% \fancypagestyle{main}{%
% 	\fancyhf{}
% 	\headstyle
% 	\footstyle
% }
%\renewcommand{\headrulewidth}{0pt}
%\renewcommand{\footrulewidth}{0.4pt}
%\renewcommand{\headrule}{\rule{\textwidth}{0.4pt}}

% appendix
\usepackage{appendix}

% -------------------------- hyperlinks & bookmarks -------------------------- %
\usepackage{aliascnt} %  \autoref
\usepackage[pdfusetitle,
bookmarks=true,
pagebackref,
colorlinks,
linkcolor=blue,
citecolor=blue,
urlcolor=blue]{hyperref}

% ------------------------------- environments ------------------------------- %
\usepackage{amsthm}
\usepackage{mdframed}
\newtheorem{dummy}{***}[section]

\theoremstyle{plain}
\newaliascnt{theorem}{dummy}
\newtheorem{theorem}[theorem]{Theorem}
\aliascntresetthe{theorem}
\providecommand*{\theoremautorefname}{Theorem}
\newaliascnt{proposition}{dummy}
\newtheorem{proposition}[proposition]{Proposition}
\aliascntresetthe{proposition}
\providecommand*{\propositionautorefname}{Proposition}
\newaliascnt{corollary}{dummy}
\newtheorem{corollary}[corollary]{Corollary}
\aliascntresetthe{corollary}
\providecommand*{\corollaryautorefname}{Corollary}
\newaliascnt{lemma}{dummy}
\newtheorem{lemma}[lemma]{Lemma}
\aliascntresetthe{lemma}
\providecommand*{\lemmaautorefname}{Lemma}
\newaliascnt{conjecture}{dummy}
\newtheorem{conjecture}[conjecture]{Conjecture}
\aliascntresetthe{conjecture}
\providecommand*{\conjectureautorefname}{Conjecture}

\theoremstyle{definition}
\newaliascnt{definition}{dummy}
\newtheorem{definition}[definition]{Definition}
\aliascntresetthe{definition}
\providecommand*{\definitionautorefname}{Definition}
\newaliascnt{example}{dummy}
\newtheorem{example}[example]{Example}
\aliascntresetthe{example}
\providecommand*{\exampleautorefname}{Example}

\theoremstyle{remark}
\newaliascnt{remark}{dummy}
\newtheorem{remark}[remark]{Remark}
\aliascntresetthe{remark}
\providecommand*{\remarkautorefname}{Remark}

\newtheorem{exercise}{Exercise}[section]
\providecommand*{\exerciseautorefname}{Exercise}

\newtheoremstyle{break}% name
  {0pt}%         Space above, empty = `usual value'
  {0pt}%         Space below
  {\itshape}% Body font
  {}%         Indent amount (empty = no indent, \parindent = para indent)
  {\bfseries}% Thm head font
  {.}%        Punctuation after thm head
  { } % HEADSPACE}% Space after thm head: \newline = linebreak
  {}%         Thm head spec

\theoremstyle{break}
\newtheorem{dummyy}{***}[chapter]
\newcommand*{\dummyyautorefname}{Dummyy}

\newaliascnt{bdefinition}{dummyy}
\newmdtheoremenv[%
  topline=true,
  rightline=true,
  leftline=true,
  bottomline=true]{bdefinition}[bdefinition]{Definition}
\aliascntresetthe{bdefinition}
\providecommand*{\bdefinitionautorefname}{Definition}

\newaliascnt{blemma}{dummyy}
\newmdtheoremenv[%
  topline=true,
  rightline=true,
  leftline=true,
  bottomline=true]{blemma}[blemma]{Lemma}
\aliascntresetthe{blemma}
\providecommand*{\blemmaautorefname}{Lemma}

\newaliascnt{btheorem}{dummyy}
\newmdtheoremenv[%
  topline=true,
  rightline=true,
  leftline=true,
  bottomline=true]{btheorem}[btheorem]{Theorem}
\aliascntresetthe{btheorem}
\providecommand*{\btheoremautorefname}{Theorem}

\newaliascnt{bsnippet}{dummyy}
\newmdtheoremenv[%
  topline=true,
  rightline=true,
  leftline=true,
  bottomline=true]{bsnippet}[bsnippet]{Snippet}
\aliascntresetthe{bsnippet}
\providecommand*{\bsnippetautorefname}{Snippet}

% ---------------------------------- autoref --------------------------------- %
\renewcommand{\sectionautorefname}{Section}
\renewcommand{\chapterautorefname}{Chapter}
%\renewcommand{\equationautorefname}{Equation}

% ------------------------------------ misc ----------------------------------- %
\allowdisplaybreaks
